\section*{Chapitre \ref{ch:b3:galois} --- Extensions de corps et
théorie de Galois}

\begin{solution}{exo:b3:galois:1}
$\sqrt3 \notin \Q(\sqrt2)$~: de $\sqrt3 = a + b\sqrt2$ ($a, b
\in \Q$), le carré donne $3 = a^2 + 2b^2 + 2ab\sqrt2$, donc $ab
= 0$ et $a^2 + 2b^2 = 3$. Si $b = 0$ alors $a^2 = 3$ impossible
dans $\Q$~; si $a = 0$ alors $2b^2 = 3$ impossible. Donc $X^2 -
3$ reste irréductible sur $\Q(\sqrt2)$, $[\Q(\sqrt2,
\sqrt3):\Q(\sqrt2)] = 2$ et $[\Q(\sqrt2, \sqrt3):\Q] = 4$ par la
tour.

Soit $\gamma = \sqrt2 + \sqrt3$. Alors $(\gamma - \sqrt2)^2 =
3$ donne $\gamma^2 - 2\gamma\sqrt2 + 2 = 3$, puis $\gamma^2 - 1
= 2\gamma\sqrt2$, $(\gamma^2 - 1)^2 = 8\gamma^2$, d'où $\gamma^4
- 10\gamma^2 + 1 = 0$. Le polynôme $X^4 - 10X^2 + 1$ est
irréductible sur $\Q$~: pas de racine rationnelle ($\pm1$
échouent)~; une factorisation $(X^2 + aX + b)(X^2 - aX + c)$
force $b + c - a^2 = -10$, $a(c-b) = 0$, $bc = 1$. Si $a = 0$
alors $b + c = -10$, $bc = 1$, discriminant $100 - 4 = 96$ non
carré dans $\Q$. Donc $[\Q(\gamma):\Q] = 4 =
[\Q(\sqrt2,\sqrt3):\Q]$, d'où égalité des corps.
\end{solution}

\begin{solution}{exo:b3:galois:2}
Les racines de $X^3 - 2$ sont $\alpha = \sqrt[3]{2}$, $\alpha
j$, $\alpha j^2$ avec $j = e^{2i\pi/3}$. L'évaluation envoyant
$\alpha$ sur $\alpha j$ définit un $\Q$-plongement $\Q(\alpha)
\to \C$ d'image $\Q(\alpha j)$, qui n'est pas $\Q(\alpha)$
(l'un est réel, l'autre non)~: $\Q(\alpha)/\Q$ n'est pas
normale. Le \omterm{thm:b3:galois:splitting}{corps de décomposition} est $L = \Q(\alpha, j)$.
$[\Q(\alpha):\Q] = 3$ et $[L:\Q(\alpha)] = 2$ ($j$ racine de
$X^2 + X + 1$, irréductible sur $\R \supseteq \Q(\alpha)$)~:
$[L:\Q] = 6$. Tout automorphisme de $\Q(\alpha)$ envoie
$\alpha$ sur une racine dans $\Q(\alpha)$, donc seulement
$\alpha$~: $\operatorname{Aut}_\Q(\Q(\alpha)) =
\{\mathrm{id}\}$, d'ordre $1 \ll 3$.
\end{solution}

\begin{solution}{exo:b3:galois:3}
$X^2 + 1$ n'a pas de racine dans $\mathbb F_3$ ($0\mapsto1$,
$1\mapsto2$, $2\mapsto2$)~: irréductible. $\mathbb F_9 =
\mathbb F_3[i]$ avec $i^2 = -1$. L'élément $1+i$~: $(1+i)^2 =
2i$, $(1+i)^4 = (2i)^2 = -4 = 2$, $(1+i)^8 = 4 = 1$, et aucun
ordre plus petit ne donne $1$ parmi les diviseurs de $8$ sauf
$8$~: générateur de $\mathbb F_9^\times$.

Sur $\mathbb F_2$~: moniques irréductibles de \omterm{def:b3:galois:extension}{degré} $1$~: $X$,
$X+1$. \omterm{def:b3:galois:extension}{Degré} $2$~: $X^2+X+1$ seul ($X^2+1=(X+1)^2$,
$X^2+X=X(X+1)$). \omterm{def:b3:galois:extension}{Degré} $3$~: $X^3+X+1$ et $X^3+X^2+1$ (sans
racine, donc irréductibles). L'identité $X^8 - X =
X(X+1)(X^3+X+1)(X^3+X^2+1)$ suit du
\cref{thm:b3:galois:finitefields} (produit des minimaux des
éléments de $\mathbb F_8$) et de l'unicité de la
factorisation.
\end{solution}

\begin{solution}{exo:b3:galois:4}
(a) $\mathbb F_7^\times$ d'ordre $6$. Un générateur a ordre
$6$~: $3^1=3$, $3^2=2$, $3^3=-1$, $3^6=1$, ordre $6$. De même
$5$~: $5^1=5$, $5^2=4$, $5^3=-1$. Pas $2$ ($2^3=1$ mod $7$).
$\mathbb F_{11}^\times$ d'ordre $10$~: $2^5 = 32 \equiv -1$,
$2^{10}=1$, ordre $10$. Aussi $6, 7, 8$ (vérifier les ordres).

(b) Le morphisme $x \mapsto x^2$ de $\mathbb F_p^\times$ a
noyau $\{\pm1\}$ (ordre $2$) et image d'indice $2$, le
sous-groupe des carrés. Or $x^{(p-1)/2} = \pm1$ pour tout $x
\neq 0$ (car $(x^{(p-1)/2})^2 = x^{p-1} = 1$), et cette valeur
vaut $1$ exactement sur le sous-groupe d'indice $2$ des
carrés (noyau de $x \mapsto x^{(p-1)/2}$). Pour $-1$~:
$(-1)^{(p-1)/2} = 1$ ssi $(p-1)/2$ est pair ssi $p \equiv 1
\pmod 4$, comme dans le~\cref{pb:b3:rings:1}.
\end{solution}

\begin{solution}{exo:b3:galois:5}
Dans une \omterm{def:b3:galois:closure}{clôture algébrique} fixée, $\mathbb F_{p^m} = \{x :
x^{p^m} = x\}$ et de même pour $n$. L'intersection est $\{x :
x^{p^m} = x = x^{p^n}\}$. Or $x^{p^a} = x$ et $x^{p^b} = x$
ssi $x^{p^{\gcd(a,b)}} = x$ (les ordres du \omterm{thm:b3:galois:finitefields}{Frobenius})~: d'où
$\mathbb F_{p^{\gcd(m,n)}}$. Le compositum $\mathbb
F_{p^m}\mathbb F_{p^n}$ est le plus petit corps contenant les
deux, donc $\mathbb F_{p^k}$ avec $m \mid k$, $n \mid k$,
minimal~: $k = \operatorname{lcm}(m,n)$.

Pour $m \mid n$, $\operatorname{Gal}(\mathbb F_{p^n}/\mathbb
F_p) = \langle\varphi\rangle$ d'ordre $n$ avec $\varphi(x) =
x^p$. Le sous-groupe $\operatorname{Gal}(\mathbb
F_{p^n}/\mathbb F_{p^m})$ est $\langle\varphi^m\rangle$,
d'ordre $n/m$.
\end{solution}

\begin{solution}{exo:b3:galois:6}
Dans $\overline{\mathbb F}_q$, $X^{q^n} - X = \prod_{a \in
\mathbb F_{q^n}}(X-a)$. Le \omterm{def:b3:galois:algebraic}{polynôme minimal} $P$ d'un $a \in
\mathbb F_{q^n}$ a \omterm{def:b3:galois:extension}{degré} $d \mid n$
(\cref{thm:b3:galois:finitefields}), et $P$ divise $X^{q^n} -
X$. Réciproquement tout monique irréductible $P$ de \omterm{def:b3:galois:extension}{degré} $d
\mid n$ a une racine dans $\mathbb F_{q^d} \subseteq \mathbb
F_{q^n}$, donc divise $X^{q^n}-X$. Les irréductibles distincts
sont premiers entre eux et le produit est \omterm{def:b3:galois:separable}{séparable}~: la
factorisation annoncée. Comparer les \omterm{def:b3:galois:extension}{degrés}~: $q^n =
\sum_{d\mid n} d\,I_d(q)$.

Conséquences~: $I_1 = q$~; $q^2 = I_1 + 2I_2$ donne $I_2 =
\frac{q^2-q}{2}$~; $q^3 = I_1 + 3I_3$ donne $I_3 =
\frac{q^3-q}{3}$~; $q^4 = I_1 + 2I_2 + 4I_4$ donne $I_4 =
\frac{q^4-q^2}{4}$. Existence~: $nI_n = q^n - \sum_{d\mid n,\,
d<n} dI_d \geq q^n - \sum_{d\leq n/2} q^d > 0$ pour $n \geq 1$.
\end{solution}

\begin{solution}{exo:b3:galois:7}
$L = \Q(\sqrt2, \sqrt3)$ est le \omterm{thm:b3:galois:splitting}{corps de décomposition} de
$(X^2-2)(X^2-3)$, \omterm{def:b3:galois:separable}{séparable}~: galoisien de \omterm{def:b3:galois:extension}{degré} $4$
(\cref{exo:b3:galois:1}). Un automorphisme envoie $\sqrt2
\mapsto \pm\sqrt2$ et $\sqrt3 \mapsto \pm\sqrt3$~: au plus $4$
choix, et $|G| = 4$ les réalise tous~: $G \cong (\Z/2\Z)^2$,
avec éléments $\mathrm{id}$, $\sigma$ ($\sqrt2 \mapsto
-\sqrt2$), $\tau$ ($\sqrt3 \mapsto -\sqrt3$), $\sigma\tau$.
Sous-groupes d'ordre $2$~: $\langle\sigma\rangle$,
$\langle\tau\rangle$, $\langle\sigma\tau\rangle$, de corps
fixes $\Q(\sqrt3)$, $\Q(\sqrt2)$, $\Q(\sqrt6)$ (noter que
$\sigma\tau$ fixe $\sqrt6 = \sqrt2\sqrt3$). Le treillis~: $\Q$
en bas, les trois quadratiques au milieu, $L$ en haut --- et
rien d'autre (\cref{thm:b3:galois:fundamental}). Pour
$(X^2-2)(X^2-3)(X^2-6)$~: le \omterm{thm:b3:galois:splitting}{corps de décomposition} est le
\emph{même} $L$ ($\sqrt6 = \sqrt2\sqrt3$), donc la réponse est
identique~: la \omterm{thm:b3:galois:fundamental}{correspondance de Galois} est un invariant de
l'\emph{extension}, non du polynôme choisi pour la présenter.
\end{solution}

\begin{solution}{exo:b3:galois:8}
(a) $G$ agit fidèlement et transitivement (irréductibilité) sur
les trois racines~: $G \hookrightarrow S_3$ avec $3 \mid |G|$~:
$G \cong A_3$ ou $S_3$. Tout $\sigma \in G$ permute les $x_i$,
et $\sigma(\delta) = \varepsilon(\sigma)\,\delta$ ($\delta$ est
alterné en les racines). Si $\Delta$ est un carré dans $\Q$~:
$\delta \in \Q^\times$ (noter $\delta \neq 0$~: \omterm{def:b3:galois:separable}{séparable}),
donc $\varepsilon(\sigma) = 1$ pour tout $\sigma$~: $G
\subseteq A_3$, d'où $= A_3$. Sinon~: $\delta \notin \Q$, donc
un certain $\sigma$ a $\varepsilon(\sigma) = -1$~: $G = S_3$.
(Dans les deux cas $\Q(\delta) = L^{G \cap A_3}$.)

(b) Avec $x_1 + x_2 + x_3 = 0$~: $u + v = 2x_1 - (x_2 + x_3) =
3x_1$. Aussi $uv = \sum_i x_i^2 + (j+j^2)\sum_{i<k}x_ix_k =
(\sum x_i)^2 - 3\sum_{i<k}x_ix_k = -3p$, en utilisant $j + j^2
= -1$ et $\sum_{i<k}x_ix_k = p$. Alors $u^3 + v^3 = (u+v)^3 -
3uv(u+v) = 27x_1^3 + 9p\cdot 3x_1 = 27(x_1^3 + px_1) = -27q$.
Donc $u^3, v^3$ sont les racines de $Y^2 + 27qY - 27p^3 = 0$
(produit $(uv)^3 = -27p^3$)~: $u^3 = \frac{-27q +
\sqrt{729q^2 + 108p^3}}{2}$, et $x_1 = \frac{u+v}{3}$ avec $v
= -3p/u$~: Cardan. La résolubilité de $S_3$ est le squelette~:
la tour $\Q \subseteq \Q(\delta) \subseteq \Q(\delta, j, u)$
adjoint d'abord une racine carrée ($\delta$, corps fixe de
$A_3$~: l'étape $S_3 \to S_3/A_3$), puis une racine cubique
($u$, car $u^3 \in \Q(\delta, j)$~: l'étape $A_3 \to \{e\}$)
--- la suite dérivée $S_3 \supset A_3 \supset \{e\}$ faite
chair.
\end{solution}

\begin{solution}{exo:b3:galois:9}
$\eta_0 + \eta_1 = \zeta_5 + \zeta_5^2 + \zeta_5^3 + \zeta_5^4
= -1$ (somme de toutes les racines $5$-ièmes de l'unité est
$0$). $\eta_0\eta_1 = (\zeta + \zeta^4)(\zeta^2 + \zeta^3) =
\zeta^3 + \zeta^4 + \zeta^6 + \zeta^7 = \zeta^3 + \zeta^4 +
\zeta + \zeta^2 = -1$ (indices mod $5$). Donc $\eta_0, \eta_1$
sont les racines de $Y^2 + Y - 1$~: $\frac{-1 \pm \sqrt5}{2}$.
Comme $\eta_0 = 2\cos\frac{2\pi}5 > 0$~: $\eta_0 =
\frac{\sqrt5 - 1}{2}$, d'où $\cos\frac{2\pi}5 =
\frac{\sqrt5 - 1}{4}$, et $\eta_1 = \frac{-1-\sqrt5}{2}$. Le
groupe $\operatorname{Gal}(\Q(\zeta_5)/\Q) \cong
(\Z/5\Z)^\times$ est cyclique d'ordre $4$~: il a un
\emph{unique} sous-groupe d'ordre $2$ ($\{\pm1\}$, i.e.\ $\zeta
\mapsto \zeta^{\pm1}$), donc $\Q(\zeta_5)$ a un unique
sous-corps quadratique (\cref{thm:b3:galois:fundamental}), qui
contient $\eta_0 \notin \Q$~: c'est $\Q(\eta_0) = \Q(\sqrt5)$.
Constructibilité~: $\cos\frac{2\pi}5$ vit dans la tour
quadratique $\Q \subseteq \Q(\sqrt5)$, et $\zeta_5$ une étape
quadratique au-dessus~: le~\cref{thm:b3:galois:wantzel}
construit le pentagone.
\end{solution}

\begin{solution}{exo:b3:galois:10}
(a) Écrire $s = S^{1/p}$, $t = T^{1/p}$. $X^p - S$ est
irréductible sur $K = \mathbb F_p(S,T)$ (\omterm{thm:b3:rings:criteria}{Eisenstein} en l'\omterm{def:b3:rings:divisibility}{élément
premier} $S$ de l'\omterm{def:b3:rings:pidufd}{UFD} $\mathbb F_p(T)[S]$, ou \omterm{def:b3:galois:extension}{degré} en $S$).
Après adjonction de $s$, $Y^p - T$ reste irréductible sur
$K(s)$ (\omterm{thm:b3:rings:criteria}{Eisenstein} en $T$)~: $[L:K] = p^2$. Tout $\alpha \in L$
s'écrit $\sum_{i,j=0}^{p-1} a_{ij} s^i t^j$ avec $a_{ij} \in
K$~; le \omterm{thm:b3:galois:finitefields}{Frobenius} $x \mapsto x^p$ (morphisme en caractéristique
$p$) donne $\alpha^p = \sum a_{ij}^p S^i T^j \in K$.

(b) Si $L = K(\gamma)$, alors $\deg\mu_\gamma = p^2$. Mais
$\gamma^p \in K$ force $\mu_\gamma$ à diviser $X^p - \gamma^p
= (X-\gamma)^p$, donc $\deg \mu_\gamma \leq p < p^2$~:
contradiction. Aucun élément primitif~: l'inséparabilité est
fatale au~\cref{thm:b3:galois:primitive}.
\end{solution}

\begin{solution}{exo:b3:galois:11}
(a) \omterm{thm:b3:rings:criteria}{Eisenstein} en $2$ ($2 \mid 4, 2$~; $4 \nmid 2$)~: $P$
irréductible. Si $\alpha$ est une racine, $[\Q(\alpha):\Q] = 5$
divise $[L:\Q] = \abs G$ ($L$ le \omterm{thm:b3:galois:splitting}{corps de décomposition})~:
Cauchy (\cref{thm:b3:groups:cauchy}) donne un élément d'ordre
$5$ dans $G \leq S_5$~; dans $S_5$, seuls les $5$-cycles ont
ordre $5$ (les ordres sont les ppcm des longueurs de cycles).

(b) $P'(x) = 5x^4 - 4$ s'annule en $\pm(4/5)^{1/4} \approx \pm
0{,}946$~: un maximum local puis un minimum local. Valeurs~:
$P(-2) = -22 < 0$, $P(0) = 2 > 0$, $P(1) = -1 < 0$, $P(2) =
26 > 0$~: trois changements de signe, et au plus trois racines
réelles (deux points critiques)~: exactement $3$ racines
réelles, donc une paire de racines complexes conjuguées.
Prendre le \omterm{thm:b3:galois:splitting}{corps de décomposition} $L$ dans $\C$~: la
conjugaison complexe envoie $L$ sur lui-même (elle permute les
racines, qui engendrent $L$) et fixe $\Q$, donc définit un
élément de $G$~; elle fixe les trois racines réelles et
échange les deux autres~: une transposition.

(c) Soit $\tau = (a\,b)$ et $\sigma$ un $5$-cycle dans $G$. Une
certaine puissance $\sigma^k$ envoie $a$ sur $b$ ($k \ne 0
\bmod 5$), et $\sigma^k$ est encore un $5$-cycle~: en
renommant, supposer $\sigma = (1\,2\,3\,4\,5)$ et $\tau =
(1\,2)$. En conjuguant, $\sigma^m\tau\sigma^{-m} =
(\sigma^m(1)\ \sigma^m(2))$~: les transpositions adjacentes
$(1\,2), (2\,3), (3\,4), (4\,5), (5\,1)$ sont toutes dans $G$~;
les transpositions adjacentes engendrent $S_5$ (toute
transposition $(i\,j)$ est produit de transpositions
adjacentes, et les transpositions engendrent). Donc $G =
S_5$, non \omterm{def:b3:groups:derived}{résoluble} (\cref{cor:b3:groups:snnotsolvable}), et
le~\cref{thm:b3:galois:solvable} conclut~: $X^5 - 4X + 2$ n'est
pas \omterm{def:b3:groups:derived}{résoluble} par radicaux.
\end{solution}

\begin{solution}{exo:b3:galois:12}
(a) $X^4 - 2$ est irréductible (\omterm{thm:b3:rings:criteria}{Eisenstein} en $2$)~:
$[\Q(\alpha):\Q] = 4$~; $\iu \notin \Q(\alpha) \subseteq \R$,
donc $[L : \Q(\alpha)] = 2$ et $[L:\Q] = 8$. L'extension est
\omterm{def:b3:galois:galois}{galoisienne} (\omterm{thm:b3:galois:splitting}{corps de décomposition} d'un \omterm{def:b3:galois:separable}{polynôme séparable}~:
les racines sont $\iu^k\alpha$), donc $\abs G = 8$. Un
automorphisme envoie $\alpha$ sur l'une des quatre racines et
$\iu$ sur $\pm\iu$~: au plus $8$ applications, toutes
réalisées. Les $\sigma$ et $\tau$ indiqués ($\sigma$ d'ordre
$4$~: $\sigma^2(\alpha) = -\alpha$, $\sigma^4 = e$~; $\tau$
d'ordre $2$) vérifient
\[
\tau\sigma\tau(\alpha) = \tau\sigma(\alpha) =
\tau(\iu\alpha) = -\iu\alpha = \sigma^{-1}(\alpha),
\qquad \tau\sigma\tau(\iu) = \iu\cdot(-1)(-1) = \iu ,
\]
plus soigneusement~: $\tau\sigma\tau(\iu) = \tau\sigma(-\iu) =
\tau(-\iu) = \iu = \sigma^{-1}(\iu)$. Donc $\tau\sigma\tau =
\sigma^{-1}$~: la présentation de $D_4$.

(b) Les dix sous-groupes de $D_4 = \langle\sigma,
\tau\rangle$~: $\{e\}$~; cinq d'ordre $2$~:
$\langle\sigma^2\rangle$, $\langle\tau\rangle$,
$\langle\sigma^2\tau\rangle$, $\langle\sigma\tau\rangle$,
$\langle\sigma^3\tau\rangle$~; trois d'ordre $4$~:
$\langle\sigma\rangle$, $\{e, \sigma^2, \tau,
\sigma^2\tau\}$, $\{e, \sigma^2, \sigma\tau,
\sigma^3\tau\}$~; et $D_4$. Corps fixes (\omterm{def:b3:galois:extension}{degré} $=$
indice)~: $\{e\} \leftrightarrow L$~; sous-groupes d'ordre $2$
$\leftrightarrow$ les cinq corps quartiques
\[
\langle\tau\rangle \leftrightarrow \Q(\alpha),\quad
\langle\sigma^2\tau\rangle \leftrightarrow \Q(\iu\alpha),
\quad
\langle\sigma^2\rangle \leftrightarrow \Q(\sqrt2, \iu),\quad
\langle\sigma\tau\rangle \leftrightarrow
\Q\bigl((1+\iu)\alpha\bigr),\quad
\langle\sigma^3\tau\rangle \leftrightarrow
\Q\bigl((1-\iu)\alpha\bigr) .
\]
Vérifications~: $\tau$ fixe le réel $\alpha$~; $\sigma^2\tau$
envoie $\alpha \mapsto -\alpha$ et $\iu \mapsto -\iu$, fixant
$\iu\alpha$~; et comme $\sigma\tau(\alpha) = \iu\alpha$,
$\sigma\tau(\iu) = -\iu$~:
\[
\sigma\tau\bigl((1+\iu)\alpha\bigr) =
(1 - \iu)\,\iu\alpha = (1 + \iu)\alpha,
\qquad
\sigma^3\tau\bigl((1-\iu)\alpha\bigr) =
(1 + \iu)(-\iu)\alpha = (1 - \iu)\alpha :
\]
chaque réflexion fixe son générateur, et le corps fixe, de
\omterm{def:b3:galois:extension}{degré} $4 = $ indice, est exactement le corps qu'il engendre
(le générateur est racine de $X^4 + 8$, irréductible).
Sous-groupes d'ordre $4$ $\leftrightarrow$ les trois corps
quadratiques~: $\langle\sigma\rangle \leftrightarrow \Q(\iu)$
($\sigma$ fixe $\iu$)~; $\{e, \sigma^2, \tau, \sigma^2\tau\}
\leftrightarrow \Q(\sqrt2)$ (les quatre fixent $\alpha^2$ à
des signes près~: $\tau(\sqrt2) = \sqrt2$, $\sigma^2(\alpha^2)
= (-\alpha)^2$)~; $\{e, \sigma^2, \sigma\tau, \sigma^3\tau\}
\leftrightarrow \Q(\iu\sqrt2)$ ($\sigma\tau(\iu\alpha^2) =
(-\iu)(\iu\alpha)^2 = \iu\alpha^2$).

(c) Galoisien sur $\Q$ $\leftrightarrow$ sous-groupes normaux
de $D_4$~: $\{e\}$, $\langle\sigma^2\rangle$ (le \omterm{ex:b3:groups:actions}{centre}), les
trois sous-groupes d'ordre $4$, et $D_4$ --- donc les corps
intermédiaires galoisiens sont $L$, $\Q(\sqrt2, \iu)$, les
trois corps quadratiques, et $\Q$. Les cinq corps quartiques
fixés par des réflexions non \omterm{def:b3:groups:normal}{normales} ne sont pas
galoisiens~: $\Q(\alpha)$ contient une racine de $X^4 - 2$
mais pas $\iu\alpha$ (il est réel) --- la conjugaison par
$\sigma$ envoie $\langle\tau\rangle$ sur
$\langle\sigma^2\tau\rangle$, exactement comme elle envoie
$\Q(\alpha)$ sur $\Q(\iu\alpha)$~: la non-normalité du
sous-groupe \emph{est} l'existence d'un corps conjugué.
\end{solution}

\begin{solution}{pb:b3:galois:1}
\textbf{1.} $\Phi_{17}$ est irréductible
(\cref{thm:b3:galois:cyclotomicirred}, ou
\cref{ex:b3:rings:eisenstein} pour un indice premier)~: $[L:\Q]
= \varphi(17) = 16$ et $G \cong (\Z/17\Z)^\times$, cyclique
d'ordre $16$ (\cref{thm:b3:galois:cyclic}). Puissances de $3$
mod $17$~:
\[
3,\ 9,\ 10,\ 13,\ 5,\ 15,\ 11,\ 16,\ 14,\ 8,\ 7,\ 4,\ 12,\ 2,\
6,\ 1
\]
--- seize valeurs distinctes~: $3$ engendre.

\textbf{2.} $G = \langle\sigma\rangle$ cyclique d'ordre $16$~;
$H_k = \langle\sigma^{2^k}\rangle$ a ordre $2^{4-k}$, et
$[H_k : H_{k+1}] = 2$. Par le théorème fondamental
(\cref{thm:b3:galois:fundamental}), $L_k = L^{H_k}$ vérifient
$[L_k : \Q] = [G : H_k] = 2^k$~: chaque $[L_{k+1}:L_k] = 2$.

\textbf{3.} $\zeta \in L = L_4$ se trouve au sommet d'une tour
d'extensions quadratiques de $\Q$~: par
le~\cref{thm:b3:galois:wantzel}, $\zeta$ est \omterm{def:b3:galois:constructible}{constructible}~; le
$17$-gone a pour sommets $\zeta^k$.

\textbf{4.} $\sigma^2$ multiplie les exposants par $9$~; les
exposants de $\eta_0$ sont les puissances paires de $3$,
\[
\{3^{2k} \bmod 17\} = \{1, 9, 13, 15, 16, 8, 4, 2\},
\]
un ensemble stable sous multiplication par $9 = 3^2$~; donc
$\eta_0$ (et de même $\eta_1$) est fixé par $H_1 =
\langle\sigma^2\rangle$~: $\eta_0, \eta_1 \in L_1$, un corps
quadratique. $\sigma$ envoie les puissances paires sur les
impaires~: il échange $\eta_0, \eta_1$. Donc $\eta_0 + \eta_1$
et $\eta_0\eta_1$ sont fixés par tout $G$~: rationnels~;
$\eta_0, \eta_1$ sont les racines d'une quadratique
rationnelle.

\textbf{5.} $\eta_0 + \eta_1 = \sum_{c=1}^{16}\zeta^c = -1$. Le
produit se développe en $64$ termes $\zeta^{a + b}$, $a$ dans
l'ensemble pair, $b$ dans l'impair. Aucun terme n'est
$\zeta^0$~: $b = -a$ est impossible, car $-1 = 16 = 3^8$ est
une puissance \emph{paire}, donc $-a$ reste dans l'ensemble
pair. Ainsi $\eta_0\eta_1 = \sum_{c \neq 0} n_c\zeta^c$ avec
$\sum n_c = 64$~; appliquer $\sigma$ fixe $\eta_0\eta_1$ (il
échange les facteurs) et permute les $\zeta^c$
transitivement sur tous les $c \neq 0$, donc tous les $n_c$
sont égaux~: $n_c = 4$ et $\eta_0\eta_1 =
4\sum_{c\neq0}\zeta^c = -4$.

\textbf{6.} $\eta_{0,1}$ résolvent $Y^2 + Y - 4 = 0$~:
$\frac{-1 \pm \sqrt{17}}2$. Numériquement, en groupant les
exposants conjugués, $\eta_0 = 2\bigl(\cos\tfrac{2\pi}{17} +
\cos\tfrac{4\pi}{17} + \cos\tfrac{8\pi}{17} +
\cos\tfrac{16\pi}{17}\bigr) \approx 1{,}56 > 0$~: $\eta_0 =
\frac{-1+\sqrt{17}}2$, $\eta_1 = \frac{-1-\sqrt{17}}2$, et
$L_1 = \Q(\eta_0) = \Q(\sqrt{17})$.

\textbf{7.} Les ensembles d'exposants~: $\beta_0$~: $\{1, 13,
16, 4\}$ = puissances $3^{4k}$~; $\beta_2$~: $\{9, 15, 8, 2\}
= 9 \times$ cet ensemble. Réunion~: l'ensemble pair~: $\beta_0
+ \beta_2 = \eta_0$~; de même $\beta_1 + \beta_3 = \eta_1$.
La multiplication par $13 = 3^4$ stabilise chaque ensemble
d'exposants de $\beta_i$~: fixé par $H_2 =
\langle\sigma^4\rangle$~; et $\sigma^2$ ($\times 9$) envoie
$\{1,13,16,4\}$ sur $\{9, 15, 8, 2\}$~: échange $\beta_0,
\beta_2$.

\textbf{8.} En développant $\beta_0\beta_2$, les seize sommes
d'exposants
\[
\{1,13,16,4\} + \{9,15,8,2\} =
\{10,16,9,3,\ 5,11,4,15,\ 8,14,7,1,\ 13,2,12,6\}
\]
couvrent $1, \dots, 16$ exactement une fois~: $\beta_0\beta_2
= \sum_{c\ne0}\zeta^c = -1$. En appliquant $\sigma$ (qui
envoie $\beta_0 \mapsto \beta_1$, $\beta_2 \mapsto \beta_3$~:
exposants $\times 3$)~: $\beta_1\beta_3 =
\sigma(\beta_0\beta_2) = -1$.

\textbf{9.} $\beta_0, \beta_2$ résolvent $Y^2 - \eta_0 Y - 1 =
0$, donc $\beta_0 = \frac{\eta_0 + \sqrt{\eta_0^2 + 4}}2$
(numériquement $\beta_0 = 2\cos\frac{2\pi}{17} +
2\cos\frac{8\pi}{17} \approx 2{,}05 > 0$, le signe $+$). De
même $\beta_1 = \frac{\eta_1 + \sqrt{\eta_1^2 + 4}}2 \approx
0{,}344$ (le contrôle numérique fixe encore le signe). $L_2 =
L_1(\beta_0)$, quadratique sur $L_1$.

\textbf{10.} $\gamma_0 + \gamma_1 = \zeta + \zeta^{16} +
\zeta^{13} + \zeta^4 = \beta_0$. Et
\[
\gamma_0\gamma_1 = (\zeta + \zeta^{16})(\zeta^{13} + \zeta^4)
= \zeta^{14} + \zeta^{5} + \zeta^{12} + \zeta^{3} = \beta_1 .
\]
Donc $\gamma_0, \gamma_1$ résolvent $Y^2 - \beta_0Y + \beta_1
= 0$~; numériquement $\gamma_0 = 2\cos\frac{2\pi}{17} \approx
1{,}865 > \gamma_1 \approx 0{,}185$~: $\gamma_0 =
\frac{\beta_0 + \sqrt{\beta_0^2 - 4\beta_1}}2$.

\textbf{11.} En chaînant~:
\[
\eta_0 = \frac{-1 + \sqrt{17}}2,
\quad
\beta_0 = \frac{\eta_0 + \sqrt{\eta_0^2 + 4}}2,
\quad
\beta_1 = \frac{\eta_1 + \sqrt{\eta_1^2 + 4}}2,
\quad
\cos\frac{2\pi}{17} = \frac{\beta_0 + \sqrt{\beta_0^2 -
4\beta_1}}4 .
\]
Numériquement~: $\sqrt{17} \approx 4{,}1231$, $\eta_0 \approx
1{,}5616$, $\eta_1 \approx -2{,}5616$, $\beta_0 \approx
2{,}0494$, $\beta_1 \approx 0{,}3441$, $\beta_0^2 - 4\beta_1
\approx 2{,}8234$, et $\cos\frac{2\pi}{17} \approx
\frac{2{,}0494 + 1{,}6803}4 \approx 0{,}93242$ --- contre
$\cos\frac{2\pi}{17} = 0{,}93247\dots$~: le petit écart est de
l'arrondi dans les affichages intermédiaires~; en portant plus
de chiffres on retrouve $0{,}932472$.

\textbf{12.} La construction demandait $[L:\Q] = p - 1$
puissance de $2$, pour qu'une chaîne complète de sous-groupes
d'indice $2$ existe. Si $p = 2^m + 1$ est premier et $m = ab$
avec $a$ impair $> 1$~: $x + 1 \mid x^a + 1$ en $x = 2^b$
montre que $2^b + 1$ divise proprement $p$ --- impossible.
Donc $m$ est une puissance de $2$~: $p = 2^{2^t} + 1$, un
\emph{premier de Fermat} ($3, 5, 17, 257, 65537$, \dots).
Réciproquement pour un tel $p$, $\varphi(p) = 2^{2^t}$ et
l'argument des questions 1--3 (ou
\cref{cor:b3:galois:impossible}) s'applique~: le $p$-gone
régulier est \omterm{def:b3:galois:constructible}{constructible} ssi $p$ est un premier de Fermat.

\textbf{13.} $\varphi(n)$ est une puissance de $2$ exactement
quand $n = 2^a p_1\cdots p_r$ avec des premiers de Fermat
distincts $p_i$ (multiplicativité de $\varphi$~; une puissance
de premier impair $p^k$, $k \geq 2$, contribue le facteur $p
\nmid 2^m$). Pour $n \leq 20$, les $n$-gones réguliers
\omterm{def:b3:galois:constructible}{constructibles} sont
\[
n = 3, 4, 5, 6, 8, 10, 12, 15, 16, 17, 20
\]
avec valeurs respectives
\[
\varphi(n) = 2,\ 2,\ 4,\ 2,\ 4,\ 4,\ 4,\ 8,\ 8,\ 16,\ 8 .
\]
Les impossibles sont $n = 7, 9, 11, 13, 14, 18, 19$, où
$\varphi(n) = 6, 6, 10, 12, 6, 6, 18$ a un facteur premier
impair.

\textbf{14.} Les carrés forment l'image du morphisme
d'élévation au carré sur le cyclique $(\Z/p\Z)^\times$,
d'indice $2$~: $\frac{p-1}2$ carrés, $\frac{p-1}2$ non-carrés,
donc les symboles somment à $0$. Alors
\[
\sum_{a=0}^{p-1}\zeta^{a^2} = 1 + 2\sum_{b \text{ carré}
\neq 0}\zeta^b
= 1 + \sum_{b\neq0}\Bigl(1 + \Bigl(\frac bp\Bigr)\Bigr)
\zeta^b = \sum_{b}\zeta^b + g = g,
\]
en utilisant $1 + \bigl(\frac bp\bigr) = \#\{a : a^2 = b\}$ et
$\sum_{b=0}^{p-1}\zeta^b = 0$. Pour $p = 17$~: les carrés mod
$17$ sont les puissances paires du générateur $3$, i.e.\ les
exposants apparaissant dans $\eta_0$ (Partie II), donc $g =
\sum_{\text{$k$ pair}}\zeta^{3^k} - \sum_{\text{$k$
impair}} \zeta^{3^k} = \eta_0 - \eta_1$.

\textbf{15.} Avec $b = c - a$ ($a, b$ parcourent les résidus
non nuls, $c = a + b$ tous les résidus)~:
\[
g^2 = \sum_c\zeta^c\sum_{a\neq0,\,a\neq c}
\Bigl(\frac{a(c-a)}p\Bigr) .
\]
Pour $c = 0$~: $\bigl(\frac{-a^2}p\bigr) =
\bigl(\frac{-1}p\bigr)$, sommé sur $p - 1$ valeurs. Pour $c
\neq 0$~: substituer $c - a = at$, i.e.\ $t = c/a - 1$~;
quand $a$ parcourt les résidus non nuls, $t$ parcourt
bijectivement les résidus $\neq -1$ (inverser~: $a = c/(1 +
t)$). Le terme devient $\bigl(\frac{a^2t}p\bigr) =
\bigl(\frac tp\bigr)$, et
\[
\sum_{t \neq -1}\Bigl(\frac tp\Bigr)
= -\Bigl(\frac{-1}p\Bigr)
\]
(la somme totale s'annule par la question 14). D'où
\[
g^2 = \Bigl(\frac{-1}p\Bigr)\Bigl[(p-1) -
\sum_{c\neq0}\zeta^c\Bigr]
= \Bigl(\frac{-1}p\Bigr)\,p = p^* ,
\]
en utilisant $\sum_{c\neq0}\zeta^c = -1$ et le critère
d'Euler $\bigl(\frac{-1}p\bigr) = (-1)^{(p-1)/2}$. Pour $p =
17$~: $(\eta_0 - \eta_1)^2 = (\eta_0 + \eta_1)^2 -
4\eta_0\eta_1 = 1 + 16 = 17$, en accord avec la Partie II.

\textbf{16.} $g^2 = p^*$ exhibe $\sqrt{p^*} = \pm g \in
\Q(\zeta_p)$, donc $\Q(\sqrt{p^*})$ est un sous-corps
quadratique. Unicité~: les sous-corps de \omterm{def:b3:galois:extension}{degré} $2$
correspondent, par la \omterm{thm:b3:galois:fundamental}{correspondance de Galois}, aux
sous-groupes d'indice $2$ du cyclique
$\operatorname{Gal}(\Q(\zeta_p)/\Q) \cong (\Z/p\Z)^\times$,
et un groupe cyclique d'ordre pair a exactement un tel
sous-groupe (les carrés). Tout corps quadratique est
$\Q(\sqrt{d})$ avec $d$ sans facteur carré, et en combinant
les corps $\Q(\sqrt{p^*})$, $\Q(\iu) \subseteq \Q(\zeta_4)$
et $\Q(\sqrt2) \subseteq \Q(\zeta_8)$ dans un
$\Q(\zeta_N)$ commun on capture tout $\sqrt d$~: le cas
quadratique de Kronecker--Weber.

\textbf{17.} Dans tout anneau commutatif, $(x + y)^q = x^q +
y^q + q(\cdots)$~: les coefficients binomiaux $\binom qk$, $0
< k < q$, sont divisibles par le premier $q$. En itérant sur
les $p - 1$ termes de $g$~:
\[
g^q \equiv \sum_{a}\Bigl(\frac ap\Bigr)^{q}\zeta^{aq}
= \sum_a\Bigl(\frac ap\Bigr)\zeta^{aq}
\pmod{q\Z[\zeta]},
\]
($q$ impair~: le symbole est inchangé). Réindexer $b = aq$~:
$a = q^{-1}b$ et $\bigl(\frac{q^{-1}b}p\bigr) =
\bigl(\frac{q}p\bigr)\bigl(\frac bp\bigr)$
(multiplicativité~; $\bigl(\frac{q^{-1}}p\bigr) =
\bigl(\frac qp\bigr)$ car le symbole d'un inverse égale le
symbole)~: $g^q \equiv \bigl(\frac qp\bigr)g$.

\textbf{18.} $g^q = g\,(g^2)^{(q-1)/2} = g\,(p^*)^{(q-1)/2}$
exactement (question 15), et le critère d'Euler dans $\Z$
donne $(p^*)^{(q-1)/2} \equiv \bigl(\frac{p^*}q\bigr) \pmod
q$, d'où mod $q\Z[\zeta]$~: $g^q \equiv
\bigl(\frac{p^*}q\bigr)g$. En comparant avec la question 17
et en multipliant par $g$~:
\[
\Bigl(\frac qp\Bigr)p^* \equiv \Bigl(\frac{p^*}q\Bigr)p^*
\pmod{q\Z[\zeta]} .
\]
Les deux membres sont des entiers rationnels~; leur
différence, $0$ ou $\pm2p^*$, vit dans $q\Z[\zeta] \cap \Z =
q\Z$ (un entier $m \in q\Z[\zeta]$ a $m/q \in \Q \cap
\Z[\zeta] = \Z$, ce dernier parce que $1, \zeta, \dots,
\zeta^{p-2}$ est une $\Q$-base dont les coordonnées
rationnelles lisent l'intégralité). Comme $q \nmid 2p^*$ ($q$
impair, $q \neq p$), la différence est $0$~:
$\bigl(\frac qp\bigr) = \bigl(\frac{p^*}q\bigr)$.

\textbf{19.} Par multiplicativité, $\bigl(\frac{p^*}q\bigr) =
\bigl(\frac{-1}q\bigr)^{(p-1)/2}\bigl(\frac pq\bigr) =
(-1)^{\frac{q-1}2\cdot\frac{p-1}2}\bigl(\frac pq\bigr)$,
donc la question 18 lit $\bigl(\frac qp\bigr)\bigl(\frac
pq\bigr) = (-1)^{\frac{p-1}2\frac{q-1}2}$~: la réciprocité.
Contrôle $(17, 3)$~: l'exposant $\frac{16}2\cdot\frac22 = 8$
est pair, donc les deux symboles doivent coïncider~; les
carrés mod $3$ sont $\{1\}$ et $17 \equiv 2$~:
$\bigl(\frac{17}3\bigr) = -1$~; les carrés mod $17$ sont
$\{1, 4, 9, 16, 8, 2, 15, 13\}$ et $3$ en est absent~:
$\bigl(\frac3{17}\bigr) = -1$. Produit $+1$, comme prédit.
Pour $x^2 \equiv 219 \pmod{383}$~: $\bigl(\frac{219}{383}
\bigr) = \bigl(\frac3{383}\bigr)\bigl(\frac{73}{383}\bigr)$.
Premier~: $383 \equiv 3 \pmod4$ et $3 \equiv 3$~: la
réciprocité donne $\bigl(\frac3{383}\bigr) =
-\bigl(\frac{383}3\bigr) = -\bigl(\frac23\bigr) = -(-1) =
+1$. Second~: $73 \equiv 1 \pmod 4$~:
$\bigl(\frac{73}{383}\bigr) = \bigl(\frac{383}{73}\bigr) =
\bigl(\frac{18}{73}\bigr) = \bigl(\frac2{73}\bigr)$ ($18 =
2\cdot3^2$), et $73 \equiv 1 \pmod 8$ fait de $2$ un carré
mod $73$ (loi supplémentaire, prouvable par $g = \zeta_8 +
\zeta_8^{-1} = \sqrt2$ dans $\Q(\zeta_8)$ par la même
méthode)~: $+1$. Total $+1$~: la congruence est \omterm{def:b3:groups:derived}{résoluble}.

\textbf{20.} Existence et unicité de la racine primitive~: si
$w$ a pour ensemble de périodes $\{d : w = u^{n/d},\ \abs u =
d\}$, le minimal tel $d_0$ divise toute autre période $d$ (si
$w$ est à la fois une $d$-puissance et une $d'$-puissance, il
est une $\gcd(d, d')$-puissance~: comparer les lettres aux
indices coïncidant modulo le pgcd, via Bézout), et le bloc de
longueur $d_0$ est primitif. En triant les $q^n$ mots par la
longueur de leur racine primitive~: $q^n = \sum_{d \mid
n}A(d)$. Comme $A$ et $d\,N_q(d)$ satisfont la même récurrence
avec les mêmes valeurs pour $n = 1$ (chacun se détermine
inductivement à partir de $q^n = \sum_{d\mid n}(\cdot)$), ils
sont égaux~: $A(d) = d\,N_q(d)$. Bijection~: un élément
$\alpha$ de \omterm{def:b3:galois:extension}{degré} $d$ donne le mot $w_\alpha$ des
coefficients de\dots\ mieux, directement~: les éléments de
\omterm{def:b3:galois:extension}{degré} $d$ dans $\overline{\mathbb F_q}$ sont les racines des
$N_q(d)$ irréductibles de \omterm{def:b3:galois:extension}{degré} $d$, chacun contribuant ses
$d$ racines distinctes (séparabilité)~: $d\,N_q(d)$ éléments
de \omterm{def:b3:galois:extension}{degré} $d$, matchant le dénombrement $q^n =
\sum_{d\mid n}\#\{\text{éléments de degré } d \text{ dans }
\mathbb F_{q^n}\}$ --- le même crible, une fois sur les mots,
une fois sur les éléments de corps.

\textbf{21.} Lemme~: $\sum_{d \mid m}\mu(d) = \mathbf
1_{m=1}$. Pour $m > 1$ fixer un premier $p \mid m$~: les
diviseurs sans carré de $m$ s'apparient en $\{d, pd\}$ avec $p
\nmid d$, et $\mu(pd) = -\mu(d)$~: la somme s'annule. Alors,
pour $f(n) = \sum_{e\mid n}g(e)$~:
\[
\sum_{d \mid n}\mu(d)\,f\bigl(\tfrac nd\bigr)
= \sum_{d \mid n}\mu(d)\!\!\sum_{e \mid n/d}\!\!g(e)
= \sum_{e \mid n}g(e)\!\!\sum_{d \mid n/e}\!\!\mu(d)
= g(n) .
\]
Avec $f(n) = q^n$ et $g(n) = nN_q(n)$
(\cref{exo:b3:galois:6})~: $N_q(n) =
\frac1n\sum_{d\mid n}\mu(d)\,q^{n/d}$.

\textbf{22.} Le terme $d = 1$ est $q^n$~; tout autre terme a
$\abs{\mu(d)q^{n/d}} \leq q^{n/2}$, et grossièrement
$\sum_{d \mid n, d > 1}q^{n/d} \leq \sum_{j \leq
n/2}q^j < 2q^{n/2}$ (géométrique, $q \geq 2$). Donc $nN_q(n)
> q^n - 2q^{n/2} \geq 0$ pour $n \geq 1$~: des irréductibles
de tout \omterm{def:b3:galois:extension}{degré} existent, et $\mathbb F_q[X]/(\pi) = \mathbb
F_{q^n}$ est (re)construit --- existence avec un recensement.
La proportion d'irréductibles parmi les $q^n$ moniques de
\omterm{def:b3:galois:extension}{degré} $n$ est $\frac1n(1 + O(q^{-n/2}))$~: le théorème des
nombres premiers de $\mathbb F_q[X]$, avec $n$ jouant le rôle
de $\log x$. Pour $q = 2$ les comptages $2, 1, 2, 3$ de
\cref{exo:b3:galois:6} matchent la formule~: p.ex.\ $N_2(4) =
\frac14(2^4 - 2^2) = 3$.

\textbf{23.} Dans le groupe abélien multiplicatif des
fractions rationnelles non nulles sur $\mathbb F_q$, poser
$F(n) = X^{q^n} - X$ et $G(n) = \prod_{\deg\pi = n}\pi$~;
\cref{exo:b3:galois:6} dit $F(n) = \prod_{d\mid n}G(d)$.
L'argument de Möbius de la question 21, écrit
multiplicativement (les exposants s'ajoutent exactement comme
les sommes), donne $G(n) = \prod_{d \mid
n}F(d)^{\mu(n/d)}$. Pour $q = 2$, $n = 2$~: $G(2) =
\frac{X^4 - X}{X^2 - X} = \frac{X(X^3 - 1)}{X(X - 1)} = X^2 +
X + 1$, l'unique irréductible quadratique sur $\mathbb F_2$,
comme il se doit.

\textbf{24.} $\omega^2 = \iu$ et $\omega^{-2} = -\iu$, donc
$g^2 = \omega^2 + 2 + \omega^{-2} = 2$. Rêve du débutant dans
l'anneau commutatif $\Z[\omega]/q\Z[\omega]$~: $g^q =
(\omega + \omega^{-1})^q \equiv \omega^q + \omega^{-q}$. La
valeur de $\omega^q + \omega^{-q}$ ne dépend que de $q \bmod
8$~: pour $q \equiv \pm1$, $\omega^q + \omega^{-q} =
\omega^{\pm1} + \omega^{\mp1} = g$~; pour $q \equiv \pm3$, en
utilisant $\omega^4 = -1$, $\omega^{3} = -\omega^{-1}$ et
$\omega^{-3} = -\omega$, donc $\omega^q + \omega^{-q} = -g$.
D'autre part $g^q = g\,(g^2)^{(q-1)/2} = g\,2^{(q-1)/2}
\equiv \bigl(\frac2q\bigr) g \pmod{q\Z[\omega]}$ par le
critère d'Euler mod $q$. En comparant et multipliant par
$g$~: $2\bigl(\frac2q\bigr) \equiv \pm2
\pmod{q\Z[\omega]}$~; si les signes différaient, $q$ diviserait
$4$ dans $\Z[\omega]$, donc dans $\Z$ ($q\Z[\omega] \cap \Z =
q\Z$~: coordonnées sur la base $1, \omega, \omega^2,
\omega^3$), impossible pour $q$ impair. Donc
$\bigl(\frac2q\bigr) = +1$ ssi $q \equiv \pm1 \pmod 8$.
Contrôle de parité~: $q = 8k \pm 1$ donne $(q^2 - 1)/8 =
2k(4k \pm 1)$, pair~; $q = 8k \pm 3$ donne $(q^2 - 1)/8 =
8k^2 \pm 6k + 1$, impair~: la formule $(-1)^{(q^2-1)/8}$
encode le scindage. Numériquement~: $3^2 = 9 \equiv 2 \pmod
7$ ($7 \equiv -1$), $6^2 = 36 \equiv 2 \pmod{17}$ ($17
\equiv 1$)~; les carrés mod $3$ sont $\{0, 1\}$ et mod $5$
sont $\{0, 1, 4\}$, aucun ne contenant $2$ ($3 \equiv 3$, $5
\equiv -3 \pmod 8$).

\textbf{25.} Tout monique $f \in \mathbb F_q[X]$ se factorise
uniquement en $\prod_\pi\pi^{e_\pi}$ sur les irréductibles
moniques~: en triant par \omterm{def:b3:galois:extension}{degré},
\[
\sum_{f \text{ monique}}t^{\deg f}
= \prod_\pi\ \sum_{e \geq 0}t^{e\deg\pi}
= \prod_\pi\bigl(1 - t^{\deg\pi}\bigr)^{-1}
= \prod_{n \geq 1}\bigl(1 - t^n\bigr)^{-N_q(n)},
\]
tous les produits $t$-adiquement légitimes (seuls les \omterm{def:b3:galois:extension}{degrés}
$\leq m$ touchent le coefficient de $t^m$, et il y a
finement d'irréductibles de chaque \omterm{def:b3:galois:extension}{degré}). Le premier membre
est $\sum_m q^mt^m = (1 - qt)^{-1}$~: l'identité. Logarithmes~:
$-\log(1 - qt) = \sum_m\frac{q^m}mt^m$, tandis que
$\sum_nN_q(n)\bigl(-\log(1 - t^n)\bigr) =
\sum_nN_q(n)\sum_k\frac{t^{nk}}k$~; le coefficient de $t^m$
donne $\frac{q^m}m = \sum_{dk = m} \frac{N_q(d)}k =
\frac1m\sum_{d \mid m}d\,N_q(d)$, i.e.\ $q^m =
\sum_{d\mid m}d\,N_q(d)$. Contrôle à la main, $q = 2$,
coefficient de $t^2$~: $N_2(1) = 2$, $N_2(2) = 1$, et $(1 -
t)^{-2}(1 - t^2)^{-1} = (1 + 2t + 3t^2 + \dots)(1 + t^2 +
\dots)$ a pour coefficient de $t^2$ $3 + 1 = 4 = 2^2$.
Enfin, la formule de la question 21 avec les diviseurs $1, 2,
3, 6$~:
\[
N_2(6) = \tfrac16\bigl(2^6 - 2^3 - 2^2 + 2\bigr)
= \tfrac{54}6 = 9,
\]
et en effet $1\cdot2 + 2\cdot1 + 3\cdot2 + 6\cdot9 = 2 + 2 +
6 + 54 = 64 = 2^6$.
\end{solution}
